% Repository technical companion, 1 October 2026.
% Retains the detailed construction methods, source families, category account
% and notation from pakemis_rewritten.tex. No inventory calculation rerun.
% This is repository technical documentation, not a designated ESSD supplement.
% Official Copernicus class remains unmodified.
\documentclass[essd]{copernicus}

% siunitx AND THE SHARED MACROS, on the same terms as paper.tex - absent until 2026-10-07,
% which meant this file did not compile at all. It uses \SI and \num throughout and no
% target built it (see the paper-v2-supplement rule in the Makefile), so the two defects
% hid each other: nothing compiled the file, so nobody saw that it could not compile.
%
% The ordering is not cosmetic and is copied from paper.tex deliberately: the class loads
% its own fontenc/inputenc, graphicx, amsmath, natbib, hyperref, tabularx, multirow, float
% and caption, and loading any of those again breaks a Copernicus build; siunitx must
% precede the macros because they call \DeclareSIUnit; and the class defines a robust
% \unit that siunitx v3 refuses to define over, so the class's copy is released first.
\let\unit\relax
\usepackage{siunitx}
% ===========================================================================
% paper_common/macros.tex - the shorthand every manuscript shares.
%
% Split from preamble.tex in the ESSD re-targeting: the national article now compiles
% under the Copernicus class, which loads its own packages and forbids most of what
% preamble.tex sets (geometry-like options, a second hyperref, placeins), but it needs
% exactly the same shorthand as the sector reports. The macros therefore live here, on
% their own, and preamble.tex inputs this file so that every other manuscript keeps
% compiling unchanged. A macro that renders differently in two documents is the defect
% this file exists to prevent; define it once.
%
% Requires: siunitx (for \DeclareSIUnit and \si) - the caller loads it.
% ===========================================================================

% Gas and pollutant names.
\newcommand{\COtwo}{CO\textsubscript{2}}
\newcommand{\COe}{CO\textsubscript{2}e}
\newcommand{\CHfour}{CH\textsubscript{4}}
\newcommand{\NtwoO}{N\textsubscript{2}O}
\newcommand{\NHthree}{NH\textsubscript{3}}
\newcommand{\NOx}{NO\textsubscript{x}}
\newcommand{\SOtwo}{SO\textsubscript{2}}
\newcommand{\PM}{PM\textsubscript{2.5}}
\newcommand{\MCF}{MCF}
\newcommand{\Mt}{\text{Mt~\COe}}
\newcommand{\tCha}{\si{\tonne\per\hectare}}
\newcommand{\AFOLU}{AFOLU}
\newcommand{\GLW}{GLW\,4}
\newcommand{\EFone}{EF\textsubscript{1}}

% siunitx: allow Mt, Tg, ha, head, yr, dB
\DeclareSIUnit{\head}{head}
\DeclareSIUnit{\yr}{yr}
\DeclareSIUnit{\dB}{dB}

% A file name is the one long unbreakable token these documents are full of -
% \texttt{waste\_air\_pollutants\_2024.csv} is wider than a two-column measure and TeX prints
% it past the margin rather than breaking it. Allowing a break AFTER each escaped underscore
% costs nothing and inserts no hyphen, which matters: a hyphenated file name is a wrong file
% name. This is one line because every such name in every document is written with \_.
\let\pkinvunderscore\_
\renewcommand{\_}{\pkinvunderscore\allowbreak}


\begin{document}
\title{Technical documentation for the national greenhouse gas emissions inventory for Pakistan, 2015--2024}
\Author[1][rehan@pakairquality.com]{Rehan}{Ahmad}
\affil[1]{Pakistan Air Quality Initiative (PAQI), Karachi, Pakistan}
\runningtitle{Technical documentation for the Pakistan greenhouse gas inventory}
\runningauthor{R. Ahmad}
\received{}
\pubdiscuss{}
\revised{}
\accepted{}
\published{}
\firstpage{1}
\maketitle

% The release cover, on the same terms as the main article and v1's supplement: the
% commit a reader needs in order to fetch exactly this release, generated from the
% release manifest rather than typed.
%% GENERATED by scripts/make_release_cover.py from outputs/release_manifest.json.
%% Do not edit: the edit will be overwritten and CI diffs the result.
\newcommand{\ReleaseCover}{%
\begin{center}\small
\setlength{\fboxsep}{4pt}%
\fbox{\begin{minipage}{\dimexpr\linewidth-2\fboxsep-2\fboxrule\relax}\footnotesize
\textbf{Release provenance.} This document describes PakEmis under the period, configuration and source boundaries specified in its methods. The identifiers below link it to the accompanying products and source register.\\[2pt]
\begin{tabular}{@{}p{0.32\linewidth}@{\hspace{0.03\linewidth}}p{0.65\linewidth}@{}}
Release version & \texttt{1.0.0rc69} \\
Commit & \texttt{d8d874622f9a} \\
Reporting period & July 2015--June 2025 (ten fiscal years) (2015-07-01 to 2025-06-30) \\
Reference & S\_TAN (household survey; TAN nitrogen) \\
FY2023--24 total & 444.61 Mt \COe{} (AR5) under S\_TAN, the reference; 449.85 Mt under L\_TAN \\
Release manifest & \texttt{ac0ffe93ad9fcd88} (SHA-256 of \texttt{outputs/release\_manifest.json}, first 16) \\
Source manifest & \texttt{c1be6d314cfce1ea} (SHA-256 of \texttt{data/manifests/source\_manifest.csv}, first 16) \\
Data-description article & \texttt{paper\_national/} (ESSD, in preparation), same release
\end{tabular}
\end{minipage}}
\end{center}}

\ReleaseCover
\begin{abstract}
This technical documentation accompanies the main article describing the PakEmis national greenhouse gas emissions inventory. It preserves the detailed construction methods, source families, reference-year category and gas account, and reporting conventions for the same inventory; it does not present a separate study or new estimates. The reporting periods are the nine fiscal years from FY2015--16 to FY2023--24, with FY2023--24 as the reference year. The documentation is intended to accompany the repository record and should be read with the main article's evaluation, limitations and data-access information. It is not designated as a journal supplementary file. Sections, equations and tables are numbered independently here.
\end{abstract}

\section{Inventory scope and source data}
\label{sec:scope}

PakEmis describes estimated anthropogenic emissions of CO$_2$, CH$_4$ and N$_2$O outside land use, land-use change and forestry (LULUCF). The reference period, FY2023--24, extends from 1 July 2023 to 30 June 2024. The historical reconstruction covers nine fiscal years beginning in FY2015--16. Ten calculation modules contribute to the account: agriculture, power, transport, residential fuel, industrial combustion, oil and gas fugitives, industrial processes and product use (IPPU), waste, railways and domestic aviation. All ten are represented in each year, but this does not imply complete national source coverage.

Modules organise the calculation; IPCC categories organise the reported inventory. These partitions differ where one module supplies several categories or a category receives contributions from several modules. In particular, public generation is reported under 1.A.1.a, while captive industrial generation contributes to 1.A.2 together with the industrial-combustion module. The transport module separates road use from a mixed off-road machinery allocation, and the railway module separates locomotive traction from a small stationary or off-road component whose finer category remains unresolved. Category and gas identifiers are therefore retained before aggregation.

The principal inputs are Pakistani administrative statistics, censuses and surveys (Table~\ref{tab:sources}). Energy records come from the National Electric Power Regulatory Authority (NEPRA), Hydrocarbon Development Institute of Pakistan (HDIP), Oil and Gas Regulatory Authority (OGRA), and Oil Companies Advisory Council (OCAC). Statistical sources include the Pakistan Bureau of Statistics (PBS), National Fertilizer Development Centre (NFDC), and Household Integrated Economic Survey (HIES). Each compiled activity record identifies its reporting period, original units, source document or dataset, transformation and applicable sector boundary. Where a source is revised, the version used in the calculation is retained rather than replaced silently. International datasets serve as inputs or comparators according to their role; the two uses are distinguished in the evaluation. The agricultural census, for example, supplies the reference livestock population, whereas FAOSTAT provides an external agricultural comparison that can share some underlying statistics and default factors.

\begin{table*}[t]
\caption{Main source families and their roles. The source register records the activity period, publication edition and location of individual records.}
\label{tab:sources}
\small
\begin{tabular}{p{0.19\textwidth}p{0.38\textwidth}p{0.35\textwidth}}
\tophline
Source family & Quantities used & Principal qualification \\
\middlehline
NEPRA and HDIP & Generation, fuel energy and coal quantities; national energy controls & Fuel-energy basis and coverage differ between aggregate and plant records. \\
OGRA and utilities & Gas sales, throughput and infrastructure & Delivered gas, feedstock, combustion and leakage have different boundaries. \\
OCAC & Petroleum sales by product and customer class & Sales categories do not always identify final combustion activity. \\
PBS and Economic Survey & Livestock, crops, population and industrial production & Census stocks, fiscal flows and industrial indices have different temporal support. \\
NFDC & Fertiliser offtake and production inputs & Offtake measures distribution rather than application to fields. \\
HIES & Household fuel quantities and generator expenditure & Survey recall, weights and the adjacent survey year affect the reconstruction. \\
Operators and producers & Railway fuel, aviation activity and industrial process quantities & Coverage and reporting basis require source-specific reconciliation. \\
International datasets & Factors, spatial proxies and external comparisons & Shared inputs and differing source boundaries limit independence. \\
\bottomhline
\end{tabular}
\end{table*}

An annual estimate does not necessarily refer to an observation made during that fiscal year. Reference livestock numbers are anchored to the 2024 agricultural census, historical stocks retain the stated calendar-year reconstruction, and household fuel estimates use the adjacent HIES 2024--25 round \citep{PBSAgri2026,PBSHIES2026}. Crop and fertiliser records follow fiscal reporting. These alignments are retained as qualifications rather than hidden by a common year label. A change between two fiscal totals can consequently reflect both recorded activity and assumptions needed to align heterogeneous sources.

Missing categories are represented through notation keys. ``Not estimated'' (NE) identifies a source without a quantified estimate; ``included elsewhere'' (IE) refers to a documented allocation to another category; ``not occurring'' (NO) and ``not applicable'' (NA) have distinct meanings. Missing cells are not replaced with numeric zeros. Material gaps include fluorinated gases, coal-mining fugitives, industrial wastewater and several industrial and combustion sources. Estimated IPPU sources include cement, ammonia, iron and steel, and lubricant use. The category--gas matrix, rather than an aggregate comparison with another inventory, defines the estimated scope.

\section{Inventory construction}
\label{sec:methods}

\subsection{Activity data and emission factors}
\label{sec:activity_factors}

We estimate emissions by source, gas and reporting period using Pakistani activity statistics and published emission factors. The reference period is 1 July 2023 to 30 June 2024, and the historical panel covers FY2015--16 to FY2023--24. Activity records retain their source period and measurement basis. An annual fuel sale, a crop harvest and a livestock census count represent different quantities even when assigned to the same reporting year. In particular, the reference agricultural herd is anchored to the 2024 census, while historical livestock estimates retain a calendar-year stock basis. Crop production and fertiliser nitrogen follow Pakistani fiscal-year records.

For a source described by an activity--factor relationship, emissions are
\begin{equation}
 E_{s,g,t}=A_{s,t}\,\mathrm{EF}_{s,g},
 \label{eq:activity_factor}
\end{equation}
where $A_{s,t}$ is the activity of source $s$ in period $t$ and $\mathrm{EF}_{s,g}$ is the factor for gas $g$. Activity and factor units are matched before aggregation. Fuel quantities are converted to energy using the corresponding net calorific value; livestock factors may instead be expressed per animal, and agricultural soil factors per unit of nitrogen. Processes with storage or transfers, including manure handling and waste disposal, are calculated from the balances described below rather than from an independent factor applied to the final category total.

The calculations use the 2006 IPCC Guidelines and relevant updates from the 2019 Refinement \citep{IPCC2006,IPCC2019}. Most categories use Tier~1 methods with country-specific activity. Selected agricultural parameters provide additional livestock characterisation, but do not constitute a uniformly Tier~2 inventory. Factors retain their process and fuel classifications: residential combustion, for example, uses the residential defaults rather than the defaults for electricity generation. Source references, units, parameter assumptions and uncertainties are retained with the inputs.

Gas masses are stored separately and converted to carbon-dioxide equivalents only for reporting. The headline AR5 factors \citep{IPCCAR5} are 1 for CO$_2$, 28 for biogenic CH$_4$, 30 for fossil CH$_4$ and 265 for N$_2$O. Methane origin is specified for each source. AR6-equivalent estimates are retained as a separate reporting basis \citep{IPCCAR6}. CO$_2$ from biomass combustion is reported as a memo item outside the greenhouse-gas total; CH$_4$ and N$_2$O from the same combustion are included. This convention does not establish that biomass use is carbon neutral: associated land-carbon changes are outside this inventory's non-LULUCF boundary.



\subsection{Electricity generation and industrial combustion}
\label{sec:energy_methods}

Power emissions are calculated from fuel-specific activity and combustion factors. NEPRA generation and fuel records provide national controls and plant information, supplemented by HDIP energy statistics \citep{NEPRA2024}. The public-generation estimate is distinguished from captive self-generation. Electricity consumed by households or industry is not added again at the point of use. Fuel properties retain their original measurement basis, including the distinction between net and gross calorific values and between energy reported in tonnes of oil equivalent and physical fuel mass.

The public-power control is a consequential unresolved input. Four readings are published for FY2023--24, each as a named configuration with its own national total. The four values are given in the main article and resolve through \texttt{outputs/national\_configurations.csv} by \texttt{configuration\_id}; they are not restated here, because a number printed in two documents is a number that can disagree with itself. The reference books the stations' own metered coal at IPCC rank-default calorific values. The plant-based alternative takes the same masses at the calorific values the plants' own generation implies, which restates the regulator's heat rates and so cannot test them; that circularity, not a preference between sources, is why the reference is selected. The legacy control is retained because a reader holding a total from an earlier release can find the account it came from, and it is not preferred: its energy fails the physical audit reported in the main article. The audit's five rescue routes and their numerical requirements are in Table~\ref{tab:coal_audit}.
%% GENERATED by scripts/make_coal_audit_table.py from outputs/pakemis_power_reference_audit.csv. Do not edit.
\begin{table*}[t]
\caption{The FY2023--24 coal supply bound and its five rescue routes, from the power reference audit. Each row states the test, its value, the audit's verdict and the change in the records that would be required to reverse it. The supply basis is imports plus production (Finance Division Statistical Supplement Table 14.3); it carries no stock term.}
\label{tab:coal_audit}
\footnotesize
\begin{tabular}{@{}>{\raggedright\arraybackslash}p{3.4cm}>{\raggedright\arraybackslash}p{2.9cm}>{\raggedright\arraybackslash}p{4.0cm}>{\raggedright\arraybackslash}p{6.1cm}@{}}
\tophline
Test & Value & Verdict & What would reverse it \\
\middlehline
B1 Tonnage required against supply & 119.5 percent of national supply & INFEASIBLE: 29531 kt required against 24719 kt supplied & supply understated by 19 percent, or energy overstated by the same \\
B2 Supply covers the same twelve months & 24719 kt & SURVIVES, with a documented caveat & supply understated by 19 percent (power alone) or 42 percent (with HDIP's other consumers) \\
B3 Inventory drawdown could supply the gap & 10335 kt of implied drawdown & SURVIVES & a drawdown of 42 percent of a full year's national supply \\
B4 Assumed lignite share & 92 percent of mass & SURVIVES & lignite share at or below 73.8 percent (inside all supply) or 41.6 percent (leaving room for other consumers) \\
B5 Fleet calorific value & 13.01 TJ/Gg & SURVIVES & a fleet net CV at or above 15.55 TJ/Gg (inside all supply) or 20.02 TJ/Gg (leaving room for other consumers) \\
B6 Fleet coverage of the station account & 99 percent & SURVIVES & not applicable in this year; coverage is 99 percent \\
\bottomhline
\end{tabular}
\end{table*}
 We retain the alternatives as separate configurations and examine the physical implications of their fuel and generation records in the main article. The module total also contains 5.55~Mt from captive generation, which is assigned to industrial combustion in the category account.

Industrial combustion combines metered gas, petroleum products, and recorded coal use in brick kilns and cement production. Gas attributed to fertiliser manufacturing is reconciled with the ammonia-process account. Petroleum customer classes are checked against those assigned to transport and power. Coal activity follows the reported consumer allocation, with separate factors and uncertainty for its use in kilns and cement plants. In the reference year, the module includes 1.564~Mt of brick-kiln coal and 4.391~Mt of cement-sector coal. The resulting industrial-combustion category is broader than metered gas alone, but remains incomplete where other fuel uses cannot be quantified.

The coal accounts require particular caution because a consumer allocation can be derived from the same balance as the national control. Agreement between the sum of consumers and supply then verifies the allocation identity rather than independently measuring each consumer. The reconstruction preserves this distinction and tests the industrial series against plant and kiln information where available. Alternative kiln activity assumptions are treated as scenarios, since a count of facilities does not specify their annual utilisation or fuel consumption.

\subsection{Road transport, railways and domestic aviation}
\label{sec:transport_methods}

Road and off-road estimates begin with petroleum customer classes and the relevant fuel properties. Retail diesel is partitioned between road activity and off-road machinery rather than assigned in full to both. The off-road allocation includes agricultural and industrial uses whose finer separation is not established by the customer class. Petrol attributed to household generators is transferred from the wider road-customer control to residential combustion. Thus, the transport-module total and the road-transport category are not interchangeable. Fuel-specific CO$_2$ factors are applied on an energy basis, while CH$_4$ and N$_2$O factors retain the appropriate mobile-combustion classification \citep{IPCC2006}.

Railway activity is based on operator fuel consumption. Volume is converted to mass and energy using documented diesel properties, and traction is separated from non-traction use. The latter is retained as a small unresolved stationary/off-road allocation rather than forced into locomotive emissions. Applying the same method to the operator's historical records permits a fiscal-year comparison with Pakistan's Biennial Update Report, discussed in the main article.

Domestic aviation is separated from international activity. Airport-related activity supports the allocation of landing and take-off emissions, whereas cruise emissions require information about flight paths that is not available at the same spatial detail. Cruise streams retain their national or administrative account and remain unallocated on the grid. International bunkers are outside the domestic national total. The aviation result should be read with its activity coverage and phase-of-flight boundary rather than as an estimate for every aircraft movement affecting Pakistan's airspace.

\subsection{Industrial processes and oil and gas fugitives}
\label{sec:process_methods}

The IPPU calculation separates process emissions from the fuels used to operate industrial facilities. Cement-process CO$_2$ uses the producer-based clinker reconstruction; a lower-clinker default case is retained separately. Cement kiln fuel is assigned to industrial combustion. Ammonia production uses the relevant gas and process-carbon account, with feedstock and process-fuel treatment reconciled to prevent a second charge to the same gas in the energy module. Iron and steel and lubricant use are additional estimated sources. The absence of an independently quantified process row does not imply zero emissions from the associated industry's energy use.

Oil and gas fugitives follow production, processing, transmission and distribution segments using the activity supported by national records. Production, utility throughput and pipeline length constrain different parts of the system. Emission factors additionally depend on operating practices, equipment and infrastructure for which Pakistani evidence is less complete. The practice and coverage alternatives are therefore reported as conditional estimates rather than assigned an unsupported probability distribution. The reference total lies at the upper end of the published practice and coverage envelope.

Gas assigned to leakage is reconciled with delivered gas used in combustion, and the boundary the deduction applies to is stated explicitly rather than implied. The deduction is taken on \emph{delivered}, post-meter gas and never on gross supply: the leak factor is itself derived downstream of the loss, so applying it to a supply-side volume would deduct a share of gas that the loss has already been removed from. The reference combustion calculation is therefore net of the identified leakage transfer for all three gases, while the leaked methane remains in fugitives, so no molecule is claimed as both fully combusted fuel and methane loss. Membership of that base is a declared list, not a residual. Of \num{9} delivered streams totalling \SI{35505}{} million~m$^3$, \num{6} carrying \SI{25211}{} million~m$^3$ (\SI{71.0}{\percent}) are inside the post-meter base --- commercial, general industry, fertiliser, captive power, cement and grid power --- and \num{3} carrying \SI{10294}{} million~m$^3$ are outside it, because residential leakage is estimated separately from an appliance-count route and vehicle CNG sits inside the road module's own control. One membership is inconsistent and is published as such: fertiliser gas is chemical feedstock, deliberately outside every combustion module, yet it is inside the post-meter base at a combustion-derived factor, and it is \SI{31.0}{\percent} of that base. The alternative accounting --- the full metered delivery treated as combusted --- is priced as a joint national configuration worth \SI{0.288}{\mega\tonne}\,CO$_2$e, and is a cross-sector adjustment no single-axis scenario can express. The ledger is \texttt{outputs/national\_gas\_closure.csv}, whose identity aborts the build if the streams stop closing on the published delivery total. None of this establishes that all national losses are measured, or that the available activity supports every fugitive segment equally; the feedstock membership and the dependence between leakage and combustion inputs remain documented limitations of the gas account.

\subsection{Agriculture and shared nitrogen accounting}
\label{sec:agriculture_methods}

Agricultural calculations cover enteric fermentation, manure management, rice cultivation, managed soils, field burning of crop residues and urea application. Reference-year livestock numbers come from the PBS Agricultural Census \citep{PBSAgri2026}; crop quantities and synthetic-nitrogen offtake come from the corresponding Economic Survey and NFDC fiscal-year publications \citep{MoF2024,NFDC2024}. The historical fiscal reconstruction uses crop and nitrogen records for each year. Enteric fermentation and manure management, together approximately 77\% of the reference agricultural estimate, retain calendar-year livestock projections outside the census anchor. A separate FAOSTAT-based calendar series provides longer context and is not substituted for the fiscal panel \citep{FAOSTAT2025}.

Enteric methane is the sum of species populations multiplied by annual per-head factors. Manure methane follows the IPCC temperature-dependent factors, while nitrogen emissions are calculated from a shared excretion ledger. For animal group $a$,
\begin{equation}
 N_{\mathrm{ex},a}=H_a r_{N,a} M_a\,365/1000,
 \label{eq:nitrogen_excretion}
\end{equation}
where $H_a$ is the animal population, $M_a$ is typical animal mass in kilograms and $r_{N,a}$ is the daily nitrogen-excretion rate in kilograms of nitrogen per 1000~kg of animal mass. The resulting nitrogen is partitioned among grazing, managed manure systems and diversion to fuel. Management shares and animal masses include declared assumptions where nationally measured parameters are unavailable. The species populations, factors, typical masses and excretion rates, with their sources, are in Table~\ref{tab:livestock_factors}. Against the First BTR's enteric CH$_4$ by species (its Figure~1), \texttt{outputs/pakemis\_enteric\_btr\_bridge.csv} shows four species --- dairy cattle, goats, sheep and camels --- agreeing to the printed precision for FY2020--21, the whole \SI{21.7}{\mega\tonne} difference sitting in buffalo (\num{65.32} against \num{48.43}) and other cattle (\num{27.45} against \num{20.80}), and two species the submission carries that this inventory does not estimate (mules and asses, horses; \SI{1.82}{\mega\tonne}).
%% GENERATED by scripts/make_livestock_factor_table.py. Do not edit.
\begin{table*}[t]
\caption{Livestock inputs, FY2023--24: population by species, the enteric CH$_4$ factor, the typical animal mass and the nitrogen-excretion rate, with the source of each. Cattle are split into dairy and other at the FAOSTAT milk-animal share (0.295); populations are the PBS 7th Agricultural Census 2024 enumeration. Factors are IPCC 2006 defaults; typical masses are Pakistan-specific values whose provenance class is stated in the parameter file.}
\label{tab:livestock_factors}
\footnotesize
\begin{tabular}{@{}lrlrrrl@{}}
\tophline
Species & Head (M) & Population basis & EF (kg CH$_4$ hd$^{-1}$ yr$^{-1}$) & $\pm$\% & Mass (kg) & N$_{ex}$ (kg N (1000 kg)$^{-1}$ d$^{-1}$) \\
\middlehline
Dairy cattle & 16.48 & census enumeration & 58 & 50 & 400 & 0.47 \\
Other cattle & 39.38 & census enumeration & 27 & 50 & 225 & 0.34 \\
Buffalo & 47.74 & census enumeration & 55 & 40 & 380 & 0.32 \\
Sheep & 44.58 & census enumeration & 5 & 50 & 35 & 1.17 \\
Goats & 95.83 & census enumeration & 5 & 50 & 32 & 1.37 \\
Camels & 1.51 & census enumeration & 46 & 50 & 480 & 0.46 \\
\bottomhline
\end{tabular}

\smallskip\noindent\footnotesize Sources. Populations: PBS 7th Agricultural Census 2024 (FY2023-24). Enteric factors: IPCC 2006 V4 Ch10 Table 10.10; IPCC 2006 V4 Ch10 Table 10.10 developing; IPCC 2006 V4 Ch10 Table 10.11 Indian Subcontinent. Typical masses: Category herd-average; Pakistan population-average adult; Population-average adult zebu; Sahiwal/crossbred adult cow. Excretion rates: IPCC 2006 V4 Ch10 Table 10.19; IPCC 2006 V4 Ch10 Table 10.19 Asia.
\end{table*}


The nitrogen available for field application is calculated after losses during manure management. These losses include volatilisation, leaching, molecular nitrogen and direct N$_2$O, using the relevant management-system factors. The pre-field volatilisation fraction, $\mathrm{FracGasMS}$, therefore reduces the manure nitrogen reaching the soil. The post-application fraction, $\mathrm{FracGASM}$, acts subsequently in the indirect N$_2$O calculation. Keeping these boundaries separate prevents one volatilisation fraction from being charged twice to the same nitrogen.

Direct soil emissions are calculated from synthetic fertiliser, applied manure, returned crop residues and nitrogen deposited during grazing:
\begin{equation}
 E_{\mathrm{N_2O,direct}}
 =\frac{44}{28}\left[
 \sum_j F_j\,\mathrm{EF}_{1,j}
 +\sum_a F_{\mathrm{PRP},a}\,\mathrm{EF}_{3,a}
 \right],
 \label{eq:managed_soil_nitrogen}
\end{equation}
where $F_j$ denotes nitrogen inputs to cultivated soils and $F_{\mathrm{PRP},a}$ denotes pasture, range and paddock deposition. Factors are expressed as N$_2$O--N per unit nitrogen; $44/28$ converts nitrogen mass to N$_2$O. The reference uses the dry-climate direct factor of 0.005 for non-rice cropland and the separate flooded-rice factor. Indirect emissions follow the volatilised and leached nitrogen pools. Alternative direct factors for cropland and grazing are evaluated separately because evidence supporting an irrigated-cropland classification does not by itself establish the appropriate factor for grazing land.

Rice methane is calculated from cultivated area, cultivation duration and the IPCC daily factor with water-regime, pre-season and organic-amendment scaling. Continuous flooding is the reference assumption; alternative water regimes are structural scenarios. Urea CO$_2$ follows urea application and its carbon content. Liming and soil-organic-nitrogen mineralisation remain not estimated where the required activity is unavailable.

Crop residues are partitioned by crop before emissions are calculated:
\begin{equation}
 R_{\mathrm{produced}}=R_{\mathrm{returned}}+R_{\mathrm{field}}
 +R_{\mathrm{fuel}}+R_{\mathrm{other}}.
 \label{eq:residue_balance}
\end{equation}
Production, residue ratios and dry-matter assumptions determine the available residue. Only the field-burning share enters agricultural combustion, after application of the combustion factor; returned residues enter soil nitrogen accounting. Fuel residues are transferred to the residential calculation, while the remaining mass retains an unresolved destination. This balance prevents the same residue from being assigned independently to field burning, soil return and household fuel.

\subsection{Residential fuel and biomass configurations}
\label{sec:residential_methods}

Residential emissions are calculated for natural gas, LPG, kerosene, coal, charcoal, firewood, dung cake, crop residues and household-generator petrol. Gas, LPG and kerosene use domestic sales statistics. Firewood, coal and charcoal use design-weighted HIES quantities \citep{HIES2019,PBSHIES2026}. Dung and residue fuel connect to the agricultural ledgers, while generator petrol is inferred from household expenditure. Purchased, own-produced, in-kind and gifted survey quantities are combined and annualised from the one-month recall period. The 2024--25 HIES round is applied to the adjacent FY2023--24 reference period; this timing mismatch remains part of the activity qualification. Historical firewood quantities interpolate per-household consumption between survey rounds using census-based household counts.

The dung calculation distinguishes the quantity diverted from agricultural management from the quantity burned in households. Cattle and buffalo fuel-diversion shares are 0.28 and 0.30, respectively. The same species populations and diversion shares determine both the removed nitrogen and the volatile solids delivered to the fuel ledger. If $V$ denotes ash-free volatile solids, total dry solids and air-dry cake are
\begin{equation}
 M_{\mathrm{dry}}=\frac{V}{1-a},\qquad
 M_{\mathrm{air}}=\frac{M_{\mathrm{dry}}}{1-w},
 \label{eq:dung_mass_bases}
\end{equation}
where $a$ is the ash fraction of dry solids and $w$ is the moisture fraction of air-dry cake. The declared fractions, 0.25 and 0.12, convert 38.77~Mt of volatile solids to 58.74~Mt of air-dry cake.

Configuration L assigns this diverted quantity to residential combustion. Configuration S instead uses the HIES household consumption estimate of 11.97~Mt of air-dry cake and retains the difference as other or unresolved use. The two are not competing estimates of one quantity, and reading them as such is what makes the gap look like an error. L is all excreted nitrogen leaving manure management for any burner, carried through equation~(\ref{eq:dung_mass_bases}); S is formed cake entering a sampled dwelling. The difference is therefore a question about destinations, and it is reported term by term --- species, age classes, moisture basis, ash fraction, preparation losses, end user and survey design, each with a comparability verdict --- in \texttt{outputs/pakemis\_dung\_balance.csv}. Three asymmetries are structural rather than numerical: the supply side covers only cattle and buffalo while the survey records cake with no species attribution, the supply side carries no preparation or handling loss at all, and the survey measures household consumption by questionnaire construction while the manure-share vocabulary names no end user, so commercial burning is outside both. Configuration L is nonetheless falsifiable, which is the useful thing to state about it: for the diversion to be household fuel at the per-household intensity the survey observes, about \SI{72.6}{\percent} of weighted households would have to burn cake against the \SI{14.79}{\percent} the survey records, or the households that do burn it would have to burn \num{4.9} times the observed daily quantity. Either is a checkable claim about Pakistan, and a national household fuel survey or a commercial cake-fuel survey would settle it. The two survey rounds differ: HIES 2018--19 gives 6.50~Mt of dung cake from 5.75~million user households (\SI{17.3}{\percent} of 33.2~million weighted households), and HIES 2024--25 gives 11.97~Mt from 5.91~million user households (\SI{14.8}{\percent} of 39.9~million). The per-user quantity, the user share and the household count all changed between rounds; the quantity moved from 6.50 to 11.97~Mt, an \SI{84}{\percent} change larger than the parameter uncertainty on either. Both satisfy the destination balance
\begin{equation}
 D_{\mathrm{diverted}}=D_{\mathrm{residential}}
 +D_{\mathrm{other/unresolved}},
 \label{eq:dung_destination}
\end{equation}
on a common mass basis. Thus, accounting closure does not resolve the disagreement between supply-derived diversion and survey-reported consumption. The configurations are reported separately, without assigning probabilities to them. Crop-residue fuel is also qualified: its agricultural ledger describes all fuel destinations, so allocating the entire fuel share to households is an upper-bound residential allocation where commercial and industrial uses are not separately measured. The nitrogen and volatile solids inside the unresolved quantity, by species and through the same chain, are in \texttt{outputs/pakemis\_dung\_residual\_ledger.csv}; the species split of the residual is an assumption that table declares.

Fuel mass and energy use consistent moisture and ash bases. As-received calorific values multiply as-received fuel mass; dry-mass pollutant factors multiply dry mass. Residential greenhouse-gas factors follow IPCC Volume~2, Chapter~2, Table~2.5. Generator petrol is estimated from reported expenditure, the administered petrol price and fuel density; the same quantity is deducted from the road-customer petrol control and assigned to residential combustion. Household electricity consumption is excluded because its generation emissions are counted in the power sector.

\subsection{Waste generation, disposal and wastewater}
\label{sec:waste_methods}

Waste calculations cover municipal solid-waste disposal, open burning and domestic wastewater. A district population panel distinguishes urban and rural residents \citep{PBSCensus2023}. It is interpolated between census observations and projected to the reference fiscal midpoint. Municipal waste generation is
\begin{equation}
 G_t=365\sum_d\left(P_{d,u,t}q_u+P_{d,r,t}q_r\right),
 \label{eq:waste_generation}
\end{equation}
where $q_u$ and $q_r$ are urban and rural daily generation rates. The adopted rates are 0.453 and 0.283~kg per person, respectively, from the Pak-EPA eight-city study \citep{PakEPA1996}. Collection is calculated using separate urban and rural shares \citep{Kaza2018}. Generation, disposal, uncollected waste, open burning and the unresolved remainder satisfy
\begin{equation}
 G_t=D_t+U_t,\qquad B_t=f_bU_t,\qquad O_t=U_t-B_t,
 \label{eq:waste_fates}
\end{equation}
with $f_b=0.60$. The uncollected, unburned remainder is retained explicitly rather than treated as a measured zero-emission fate. Alternative unmanaged-disposal assumptions are evaluated as scenarios.

Solid-waste methane is calculated by first-order decay over the deposit history. A worked annual example reproducing the engine's FY2023--24 value from the deposit series, the 1960 spin-up and the stated parameters is in \texttt{docs/waste\_fod\_worked\_example.md}. Let $S_t$ be the accumulated decomposable degradable organic carbon, $C_t^{\mathrm{in}}$ its addition from disposed waste and $C_t^{\mathrm{dec}}$ the amount decomposed during period $t$. The carbon balance and methane conversion are
\begin{align}
 S_t&=S_{t-1}+C_t^{\mathrm{in}}-C_t^{\mathrm{dec}},\\
 E_{\mathrm{CH_4},t}&=C_t^{\mathrm{dec}}F\frac{16}{12},
 \label{eq:waste_fod}
\end{align}
with $C_t^{\mathrm{in}}=W_t\,\mathrm{DOC}\,\mathrm{DOC}_f\,\mathrm{MCF}$ for disposed wet waste $W_t$. DOC is degradable organic carbon per unit waste, $\mathrm{DOC}_f$ is its decomposable fraction, MCF represents disposal conditions and $F$ is the methane fraction of landfill gas. Decomposition follows the deposit history with $k=0.065$~yr$^{-1}$. The reference uses $\mathrm{DOC}_f=0.5$, $\mathrm{MCF}=0.6$ and $F=0.5$, with no methane recovery or cover oxidation credited. The history extends to 1960 using a declared pre-panel assumption; a no-spin-up scenario assesses its influence. Each historical year is obtained by reconstructing deposits and running the decay model to that year's endpoint. Generation rates, collection, composition and sanitation shares remain fixed, so the series represents demographic change and accumulated deposits rather than observed changes in waste-management practice.

Open-burning factors are applied to $B_t$ on their stated wet- or dry-mass bases. Fossil carbon contributes reportable CO$_2$; biogenic CO$_2$ is a memo item. CH$_4$ and N$_2$O are included for both carbon origins. Domestic wastewater methane follows population-weighted sanitation pathways:
\begin{equation}
 E_{\mathrm{CH_4}}=\sum_j \mathrm{TOW}_j B_0\mathrm{MCF}_j,
 \label{eq:wastewater_methane}
\end{equation}
where TOW is degradable organic load, calculated from population, per-capita biochemical oxygen demand and the industrial--commercial correction. Sanitation shares follow the WHO/UNICEF Joint Monitoring Programme \citep{JMP2023}. The reference uses 40~g BOD per person per day, a correction of 1.25 and $B_0=0.6$~kg CH$_4$ per kilogram BOD. Sewer MCF values of 0.5 and 0.1 represent alternative stagnant-sewer and receiving-water interpretations. Wastewater N$_2$O is calculated separately from protein-derived nitrogen and the effluent factor; it is therefore unchanged in scenarios varying only BOD or sewer MCF.

The two gases do not share a denominator, and the asymmetry is deliberate but consequential enough to state. The methane side splits the population across five sanitation pathways with a methane correction factor for each, and open defecation carries an explicit zero rather than being dropped, because material left on open ground decomposes aerobically and the Guidelines offer no anaerobic pathway for it. The nitrous-oxide side applies the effluent factor to nitrogen \emph{generated by the whole population}, with no pathway split and the sludge-nitrogen term at its default of zero, following the 2006 Tier~1 effluent equation. This is faithful to the method and is not a residual. The absent sludge-nitrogen term is a DEFAULT ZERO rather than a demonstration, and the distinction matters: we hold no measurement of sludge removal in Pakistan, and the grounds for setting the term to zero are that centralised treatment is nearly absent and that the sludge-nitrogen term of the Tier~1 equation (IPCC 2006, Vol.~5, Ch.~6, Eq.~6.8) takes its default value of zero. Those grounds make zero the defensible default; they do not establish that no nitrogen leaves the stream as sludge. What would settle it is a sludge-production measurement at the treatment capacity that does exist, which would enter as a deduction from the denominator and could only reduce this row. It does, however, make the N$_2$O denominator a hypothesis about where effluent nitrogen reaches a nitrifying environment rather than an observation. A strict aquatic-discharge reading --- restricting the denominator to the sewer-connected population, \SI{74.0}{} of \SI{247.2}{} million people --- would divide the N$_2$O row by \num{3.34}. That alternative is not adopted, because the Guidelines' own Tier~1 denominator is the whole population and substituting a narrower one would be a method change rather than a correction; it is recorded so the scope of the assumption is visible beside the factor's own width.

Industrial wastewater and other unquantified waste categories retain explicit not-estimated status. The structural alternatives the module prices one at a time --- the sewer mixture endpoints, the BOD rate, no spin-up, the HIES route shares and the unresolved-fate members --- are tabulated with their conditional intervals in \texttt{outputs/waste\_uncertainty\_scenarios\_2024.csv}; the historical generation, collection and sanitation series enter through the deposit history and the pathway shares those rows vary.


\subsection{Monthly allocation and temporal evidence}
\label{sec:temporal}

Annual emissions are allocated to months using source-specific profiles. For stream $s$, gas $g$ and fiscal year $y$,
\begin{equation}
 E_{s,g,y,m}=E_{s,g,y}\,p_{s,y,m},\qquad
 \sum_{m=1}^{12}p_{s,y,m}=1.
 \label{eq:monthly}
\end{equation}
The profile is independent of the gas reporting metric: the same activity timing is applied before conversion to CO$_2$ equivalents. Fiscal months run from July to June, and calendar-year sources retain the transformation used to align them. This approach is consistent with the role of temporal profiles as allocations of annual emissions \citep{Crippa2020}; it does not turn monthly proxies into direct emission observations.

Four evidence classes distinguish the origin of a profile. TE1 uses observed source-month activity; TE2 uses reconstructed monthly administrative activity; TE3 uses a process or activity-shape proxy; and TE4 uses a weakly constrained or flat allocation. A flat profile represents a constant daily rate, so monthly mass is proportional to elapsed days, including leap years. It is not twelve equal monthly masses and should not be interpreted as evidence that the process lacks seasonality.

Plant dispatch, petroleum sales, fertiliser offtake, industrial production and fire detections contribute where the stream and reporting period support their use. Their meanings are retained: electricity generation need not identify the monthly fuel mix, fertiliser sales need not coincide with field application, and a production profile is not automatically a profile for downstream deliveries. Historical records are used for their own years where available. Borrowed profiles retain the source year, and unsupported historical detail is not inferred from the reference year's seasonal shape without qualification.

Monthly inputs are checked for date order, missing values, unit consistency and agreement with annual source totals. Reconstructed series are compared with independently printed tables where available. These checks address extraction and timing; they do not independently validate the timing of emissions. Profile alternatives are evaluated separately from changes to annual activity so that the effects on national magnitude and within-year distribution remain distinguishable.

\subsection{Spatial allocation and geographical support}
\label{sec:spatial}

Spatial products use a fixed 0.05$^{\circ}$ geographical lattice and administrative crosswalks. Point sources use supported facility locations; distributed sources use administrative activity or spatial proxies appropriate to the process. Agricultural crop panels, livestock distributions, settlement information, industrial facilities and transport infrastructure therefore have different roles. The reference administrative frame contains 152 second-level units based on the stated geoBoundaries 2019 frame. At the first level the inventory carries five units, and the fifth needs stating because it is a residual rather than a province: Punjab, Sindh, Khyber Pakhtunkhwa and Balochistan are reported individually, and \emph{Other} aggregates Gilgit-Baltistan, Azad Kashmir and Islamabad Capital Territory. Emissions arising in those three territories are therefore inside the national total and inside the national denominator, and are reported together rather than separately --- the activity statistics that drive the provincial controls generally print four provinces and a remainder, so a finer split would be an allocation assumption and not a reported figure. Where a source does distinguish them, such as the curated station record, the distinction is preserved in the underlying evidence and collapsed only at the reported first level.
Table~\ref{tab:other_units} fixes the vocabulary used for these three territories throughout this document and the data products.
\begin{table}[h]
\caption{The three territories aggregated as \emph{Other}: where each sits in the frame, the total and the reporting. Source: \texttt{src/pkinv/core/geography.py} (\texttt{ADM1\_TO\_PROVINCE}) and \texttt{data/shared/gis/stream\_geometry.csv}.}
\label{tab:other_units}
\begin{tabular}{llll}
\tophline
Unit & In the declared boundary frame & In the national total & Reported at the first level as \\
\middlehline
Gilgit-Baltistan & yes & yes & Other \\
Azad Jammu \& Kashmir & yes & yes & Other \\
Islamabad Capital Territory & yes & yes & Other \\
\bottomhline
\end{tabular}
\end{table}
``Inside the frame'' means the unit is part of the geoBoundaries ADM1 geometry the grid is cut to; ``in the national total'' means its emissions are inside the reference total and its denominator; ``outside the statistical panels'' (the phrase used in the map legend) means the provincial activity statistics that drive the controls print the four provinces and a remainder, so these units are allocated from the remainder rather than from a line of their own. Historical allocations on this frame describe estimates within a common geography rather than observations on each year's changing district boundaries.

For a placed stream, spatial weights satisfy
\begin{equation}
 E_{s,g,y,i}=E_{s,g,y}\,w_{s,y,i},\qquad
 \sum_i w_{s,y,i}=1,
 \label{eq:spatial}
\end{equation}
where $i$ denotes the spatial unit. Mass that cannot be placed at the requested resolution is carried in an explicit unallocated term. National conservation is assessed using both parts:
\begin{equation}
 \sum_i E_{g,y,i}+E_{g,y,\mathrm{unallocated}}
 =E_{g,y,\mathrm{national}}.
 \label{eq:spatial_closure}
\end{equation}
Raster, administrative and physical-source views require their declared crosswalks and any outside-boundary terms. They are not separate emissions to be added together, and an independent sector raster should not be assumed to reproduce a district table built with another allocator.

Each historical year uses that year's annual emission streams. Spatial weights are drawn from the same year when supported, while timeless locations and borrowed or substituted patterns remain identified. A change in a grid cell can therefore arise from annual activity, a changing spatial input or a fixed pattern applied to a new total. Users should inspect the weight provenance before interpreting local trends.

The closure identities between the products --- monthly sums equal to the annual estimate, and placed plus unallocated mass equal to the national stream ledger --- hold per gas to the printed precision (Table~\ref{tab:closure}).
%% GENERATED by scripts/make_closure_table.py from construction_quantities.csv. Do not edit.
\begin{table}[t]
\caption{Closure of the spatial products for 2023--24, per gas, in Gg of the native gas: the mass placed on the grid plus the mass retained unallocated equals the national stream ledger. The reference total is 451.86 Mt CO2e (AR5).}
\label{tab:closure}
\footnotesize
\begin{tabular}{@{}lrrr@{}}
\tophline
Gas & Placed (Gg) & Unallocated (Gg) & Ledger (Gg) \\
\middlehline
CO$_2$ & 160694.749708 & 249.458356 & 160944.208064 \\
CH$_4$ & 8827.313149 & 0.001741 & 8827.314890 \\
N$_2$O & 158.969280 & 0.006963 & 158.976243 \\
\bottomhline
\end{tabular}
\end{table}


The allocation inputs are versioned, licensed and hashed, and the record of them is machine-readable rather than descriptive. Five committed 0.05$^{\circ}$ proxy rasters carry the distributed streams: WorldPop constrained population (2020), ESA WorldCover v200 built-up and cropland classes (2021), FAO Gridded Livestock of the World~4 weighted by species emission factors (2020), and a derived cropland~$\times$~livestock manure-application surface. Two further surfaces are manual acquisitions held in the repository: MapSPAM 2020 rice physical area (V2r2) and VIIRS fire detections for residue burning. Each carries its licence and a SHA-256 digest, and each digest enters the release manifest, so a re-derivation can establish identity without the manifest. \textbf{Every surface predates the reference reporting period: built-up and cropland observe 2021 and population 2020, a declared vintage offset of two to three years against the activity they allocate.} A spatial proxy here is a versioned snapshot, never a measurement of the reporting period. A separate 2023 30~m classified-paddy raster is used only as a rice crosscheck instrument and is never an allocation surface, and its result is reported rather than merely cited. The instrument is GCD-Rice 2023 (30~m classified paddy for calendar year 2023, \SI{88.9}{\percent} reported classification accuracy, which is an accuracy of the paddy map and not of emission location), independent of MapSPAM in both its observation and its method; it observes calendar 2023 against a FY2023--24 reporting period rather than preceding it by several years. Compared as within-province ADM2 shares --- the quantity the allocation actually uses --- the primary statistic is the threshold-free one: half the L1 distance between the two surfaces' district shares, per province, given below. At a declared reading line of \num{2.0} percentage points of a unit's own provincial rice mass, the two surfaces agree in \num{76} of \num{126} districts and disagree in \num{50}. Aggregated per province, the fraction of provincial rice mass the own-year instrument would place in different districts is \SI{25.3}{\percent} in Punjab, \SI{38.0}{\percent} in Sindh, \SI{62.3}{\percent} in Khyber Pakhtunkhwa, \SI{34.3}{\percent} in Balochistan and \SI{41.5}{\percent} in the residual province. These are disagreements about \emph{where} rice emissions sit inside a province, not about how much there is: provincial and national totals are controlled elsewhere and no total moves. Coverage is \num{126} of the canonical \num{152} ADM2 units --- the instrument is rice-growing area, so districts with no detected paddy on either surface carry no comparison --- and the full table is \texttt{outputs/pakemis\_rice\_geography\_crosscheck.csv}. We report this as the size of an open question about the rice surface, not as a correction: adopting the own-year instrument would be a declared re-base of the agriculture geography and is not made here. Two products added in v3 make the evidence behind the district geography readable without the registry: \texttt{outputs/pakemis\_spatial\_basis\_shares.csv} gives the mass-weighted share of each year under each evidence class, with the classes defined in plain words in \texttt{pakemis\_spatial\_basis\_classes.csv} and the same shares per reporting sector in \texttt{pakemis\_spatial\_basis\_shares\_by\_sector.csv} (FY2023--24: \SI{61.5}{\percent} on a timeless structural surface, \SI{30.9}{\percent} on that year's observed activity, \SI{7.5}{\percent} on a declared stand-in surface, \SI{0.06}{\percent} retained at province level); and \texttt{outputs/pakemis\_power\_fixed\_pattern\_exposure.csv}, with its station table \texttt{pakemis\_power\_station\_status.csv}, prints per period and province the coal mass the reference year's provincial vector books against what the period's own station pool would book, beside the Global Energy Monitor commissioning window of every station (no station carries coal outside its window; the fixed provincial vector differs from the period's own pool by up to \SI{11.6}{\mega\tonne} in one province-year, and by at most \SI{0.17}{\mega\tonne} in the reference year).

Four registers make the allocation inspectable at the grain at which it is decided. A per-stream proxy table states which surface each stream uses; a geometry table and a geometry-set table state the shapes it is placed within; and the per-stream-year weight registry publishes 1818 rows and 27 columns, including the geometry and surface vintages separately, the share of mass placed by a substituted surface, the period a borrowed pattern came from, and a pattern hash that makes ``borrowed'' a checkable property rather than a label.

That registry also carries a four-class vocabulary for the time basis of each placement, and the classes are deliberately narrow. \emph{Annual observed activity} (795 stream-years) requires both that the evidence has a time axis and that the engine reads it for that period. \emph{Static structural} (615) is a timeless surface --- a raster snapshot or a declared judgement --- in any period, \emph{including the reference year}: being the reference year does not make a 2021 snapshot an observation of FY2023--24, and an earlier vocabulary that implied otherwise labelled 196 stream-years and \SI{632.9}{\mega\tonne} in a way that invited exactly that reading. \emph{Reference shape backcast} (306) is the reference year's pattern rescaled to another year's control, and \emph{substituted surface} (46) is a placement where the declared surface had no evidence. \num{27} stream-years are never placed below the first administrative level and are labelled as such, and \num{29} carry no mass at all in their period --- a stream-year with nothing to place has no geography to grade --- so they are held as explicit rows rather than omitted. Those six values account for all \num{1818} rows.

Substitutions are fail-closed rather than silent. Every fallback the cube takes is written to a declared-fallback product, and a release may ship only fallback \emph{classes} that appear in a committed allow-list with a written justification. A class that appears in the fallbacks product and not in the allow-list --- a deleted proxy cache, a renamed raster, a new stream with no registry row --- refuses the release fixpoint. When a manual acquisition landed, its former absence class was deleted from the allow-list, so a checkout without that file now refuses to build rather than quietly substituting.

The gridded variable is mass per cell and period, expressed in native-gas units, rather than a flux density. Cell area and the geographical fraction used in allocation accompany the product. Conversion to a flux requires the applicable area and period duration; the fraction must not be applied a second time if it is already included in the allocated mass. For an annual cell mass $E_i$ in gigagrams, a full-cell average flux is $10^6E_i/(A_i\Delta t)$ in kg~m$^{-2}$~s$^{-1}$, where $A_i$ is cell area in square metres and $\Delta t$ is the period in seconds. An in-country-area convention instead uses the corresponding covered area and must be labelled accordingly.

\paragraph{Provenance of the gridded figures.}
The paragraphs below are generated by the figure script itself
(\texttt{scripts/make\_fig\_ghg\_grid.py} $\rightarrow$
\texttt{paper\_national/figures/fig\_ghg\_grid\_provenance.tex}) from the same arrays the
figures plot, so they cannot describe a different build. They were previously typeset inside
the figure at 6.3\,pt, which is what forced a page larger than either Letter or A4; as prose
they reflow, are searchable with the rest of the text, and carry the same numbers.
%% GENERATED by scripts/make_fig_ghg_grid.py. Do not edit: the edit is overwritten.
%% The provenance of the gridded figures, as prose.
\newcommand{\GridProvenance}{%
\par\smallskip\noindent\textbf{WHAT THE MAPS ARE.} The map figure draws the released 0.05-degree annual grid for 2023--24 (release 1.0.0rc68, the reference configuration), read and not recomputed. Each cell carries the mass emitted in it over the fiscal year, in Gg of the native gas: no GWP is applied, no CO2e is formed, and the value is not divided by cell area, so it is a mass and not a flux density. The cube's three spatial views (physical, origin, accounting) are disjoint and partition the mass, so a cell's value is their sum.
\par\smallskip\noindent\textbf{PERIOD.} 2023--24 is a Pakistani fiscal year, 1 July 2023 to 30 June 2024. It is not a calendar year and is not relabelled as one. The evidence figure runs the nine fiscal years 2015--16 to 2024--25.
\par\smallskip\noindent\textbf{SOURCES INCLUDED.} Every estimated source stream of the 10 calculation modules: agriculture, aviation, fugitives, industry, ippu, power, railways, residential, transport, waste. F-gases, lime (2.A.2), LULUCF and commercial and institutional buildings are not estimated by this release and are therefore absent from the maps. Their absence is a coverage gap, not a zero.
\par\smallskip\noindent\textbf{THE MASS NOT ON THE MAPS.} 606.69 Gg CO2 (0.60669 Mt; 0.38\% of national CO2), with 0.021748 Gg CH4 and 0.144838 Gg N2O (the three gases together 0.64573 Mt CO2e on AR5), is domestic-aviation cruise, which has no native location below the province and is held at province level. It is the fifth class of the evidence figure and is not spread to make the maps close. The identity grid + unallocated = national stream ledger holds for all 30 period-by-gas combinations; the worst residual is 0.004 ppm of the ledger, which is float32 storage and not a disagreement.
\par\smallskip\noindent\textbf{PROJECTION, SCALE AND THE THREE STATES OF A CELL.} EPSG:4326 (WGS 84), drawn plate carree with the aspect set for latitude 30.4 degrees N; no reprojection and no interpolation, so each native cell is drawn as itself. The lattice is equal-angle: a cell's area runs from 24.678 km2 at the northern edge to 28.298 km2 at the southern, which changes the apparent size of a cell and not its colour, because the quantity is a mass. The colour scale is logarithmic and separate for each gas, so equal colours on two maps are not equal masses. Its floor is the decade at or below the smallest of the cells that together hold 99.9 per cent of the gas's placed mass, and its ceiling the next of 1, 2 or 5 times a power of ten above the largest cell. Of the 33,674 valid cells, those holding exactly zero mass (CO2 0; CH4 0; N2O 3,432) are drawn in a neutral grey that is not on the scale, and those below the floor (CO2 13,229; CH4 4,301; N2O 3,048), which together hold under 0.1 per cent of the mass, take one pale colour marked by the pointed end of the bar. A cell cut by the coast or the border holds only the mass of its in-country part. Cells outside every district carry no value and are left blank.
\par\smallskip\noindent\textbf{BOUNDARY FRAME.} geoBoundaries gbOpen PAK ADM1/ADM2, represented year 2019, release commit 9469f09 (Runfola et al. (2020), PLOS ONE 15(4), e0231866; Public Domain). The gridded frame includes the 11 territory units of Gilgit-Baltistan and Azad Jammu and Kashmir, drawn dashed: they are inside the gridded frame and outside the national statistical tables, and several comparator products draw them differently, so a cell-by-cell comparison against those products is not like-for-like there. Districts carry the 2023 census vocabulary on 2019 geometry.
\par\smallskip\noindent\textbf{WHAT SUPPORTS THE PLACEMENT.} No cell is an observation. In the reference year 68\% of CO2, 93\% of CH4 and 29\% of N2O are placed by a weight that is not that year's own observed activity; in 2015--16 those shares are 68\%, 93\% and 30\%. The surfaces in use in the reference year are: ESA WorldCover 2021; ESA WorldCover 2021 + OCAC area pages (per period); GLW 2020; MapSPAM 2020; VIIRS active-fire detections (per period); World Bank satellite flare-site estimates; source calendar year recorded separately; WorldPop 2020; derived: ESA WorldCover 2021 cropland x GLW 2020. A surface's vintage is not a measure of its error, and not every surface predates the year it places: none is later than the reference year. Per-cell uncertainty is deliberately not stored: within a category the parametric error is perfectly correlated, so a per-cell standard deviation would invite addition in quadrature and understate every aggregate.
\par\smallskip\noindent\textbf{CONCENTRATION.} The concentration curves rank all 33,674 valid cells of the reference year from the highest mass, zero-valued cells included. Half of the placed mass sits in 64 cells for CO2, 3,055 for CH4 and 3,516 for N2O. These counts belong to this grid, this mask and this allocation: a finer display grid would add no evidence about where sources are.
}

\GridProvenance{}

\subsection{Uncertainty, dependence and structural alternatives}
\label{sec:uncertainty_methods}

The inventory combines measured activity, emission factors, conversion factors and allocation assumptions. These inputs affect different properties of the estimate. An uncertain livestock emission factor changes annual national emissions, whereas uncertain provincial shares change their location without changing the national mass. We therefore distinguish parameter uncertainty from structural choices, temporal profiles, spatial allocation and incomplete source coverage. The reported probability intervals describe parameter uncertainty conditional on a specified inventory configuration. Alternatives for which probabilities are unsupported are reported separately.

The parameter registry connects each emission stream to its activity, conversions and emission factors. A stream can depend on several parameters, and a parameter can enter several streams. For example, a power-sector gas stream depends on its activity, the conversion from tonnes of oil equivalent to energy, and the carbon content of natural gas. These terms are sampled separately, with the same draw used wherever a shared parameter enters the calculation. This representation preserves common inputs without imposing the full uncertainty of one stream on another. Remaining uncertainty that is not assigned to a named parameter is retained as a stream-specific residual.

We report both IPCC Approach~1 error propagation and Approach~2 Monte Carlo results. Approach~1 combines the declared component uncertainties, preserving the dependence specified within its groups. Approach~2 propagates 100\,000 realizations through the national calculation using national driver seed~42, with the inherited sector seeds retained. Waste and fugitives contribute the persisted draws from their own sector engines. Waste is recalculated with its material-fate identities preserved in every realization; fugitives retains its bounded split-triangular input distributions and shared activity draws. The other modules use mean-preserving lognormal multipliers for their sampled parameters. The national ensemble consequently preserves these different sector distributions rather than fitting one distribution to the national total. Source files, seeds and checksums identify the draws used.

For configuration $s$, the national total in realization $d$ is
\begin{equation}
 T_y^{(d)}(s)=\sum_i \mathrm{GWP}_{g(i)}\,
 E_{i,y}^{(d)}(s),
 \label{eq:uncertainty_total}
\end{equation}
where $E_{i,y}^{(d)}$ is the native-gas emission of stream $i$ in year $y$, and $g(i)$ identifies its gas. Percentiles are calculated after aggregation. Thus, a national interval retains the dependence present in the draws; it is not obtained by summing sector percentile limits. We distinguish the deterministic reference total, used for accounting comparisons, from the simulated median, used to summarize the conditional distribution.

Two dependency settings examine the treatment of fuel-carbon factors. In the first, sharing stops at the module boundary. In the second, all eligible combustion streams using the same physical fuel-carbon factor share its draw across modules. Separate groups represent diesel, gasoline, natural gas, residual fuel oil, coal, jet kerosene, other kerosene and LPG. Process emissions and other non-combustion carbon sources are excluded from these groups. Activity and conversion parameters retain their own dependencies. Common random numbers support comparison between settings, and batch standard errors quantify simulation variability. These settings test specified dependency assumptions; they do not estimate the true covariance of fuel properties or establish bounds on every possible dependency structure.

Parameter importance is assessed by variance decomposition, and what is decomposed is \textbf{a surrogate of the inventory, not the inventory}. The sector engines are evaluated once to obtain the component list and its reference values; the sampled model is then a closed-form mean-preserving multiplicative lognormal over those components, in which each component is scaled by the dependency groups it belongs to and by a stream-specific residual. Nothing is re-run through the engines, so for the two modules whose ensembles are injected rather than re-expressed --- waste and fugitives --- \textbf{the indices are not a decomposition of the reported ensemble and must not be read as one}.

The size of that difference is measurable and is largest exactly where it matters. The wastewater effluent factor is sampled in the waste engine as a bounded split-triangular distribution in log space, which is \emph{median}-preserving: median \num{0.005}, mean \num{0.01608}. The surrogate matches its log-variance exactly but is \emph{mean}-preserving: mean \num{0.005}, median \num{0.002196}. The variance of the factor is therefore \num{7.606e-4} as sampled against \num{1.046e-4} in the surrogate, a factor of \num{7.28}. Its total-effect index is understated by a similar factor. Measured by Jansen pick-freeze on the waste engine's own draws --- resampling this factor alone while every other parameter holds its realisation, at \num{10000} draws per pass --- and divided by the \emph{same national denominator} the published decomposition uses, the index is \num{0.2678} against \num{0.04895} as published, a factor of \num{5.47}. The denominator has to be stated because it decides the number: the same numerator against the waste module's own variance gives \num{0.8296}, and quoting that beside a national index would be an artefact of two denominators rather than a finding. The \num{0.2678} is the waste engine's own numerator over the surrogate's national denominator. It is not an index of the published ensemble and it is not placed in the published ranking, whose entries are surrogate numerators over the same surrogate denominator; the two numbers share a denominator and nothing else. The two analyses answer different questions and should not be read as one.

The same factor supplies the inventory's largest \emph{parameter-ceiling} scenario. It is not a structural endpoint: nothing about the system's structure changes, and the quantity priced is the ceiling of a parameter's own sampled support, so it belongs with parameter uncertainty and not beside the configuration alternatives. Pricing the published range's ceiling adds \SI{148.6}{\mega\tonne}\,CO$_2$e to the national total, which exceeds the reference ensemble's \SI{97.5}{\percent} point; it does not follow that the endpoint lies outside what the ensemble samples. The ceiling itself is the \emph{supremum} of the sampled support --- the split-triangular density falls to zero there, so \num{0.25} is attained with probability exactly zero and the realised maximum over \num{100000} draws is \num{0.246007} --- while the endpoint's \emph{emission level} --- a national total of \SI{600.501}{\mega\tonne}\,CO$_2$e --- is reached by the ensemble in \num{12} of the \num{10000} persisted draws. That is an empirical Monte Carlo estimate of \SI{0.12}{\percent}, not an exact probability: twelve events carry an exact binomial \SI{95}{\percent} interval of \SIrange{0.062}{0.210}{\percent}, so the frequency is known only to within about a factor of three, and the point estimate corresponds to roughly one draw in \num{833}; and the ensemble's \SI{99.9}{\percent} point, \SI{602.779}{\mega\tonne}, sits above it. The endpoint is therefore a stress test the ensemble already samples and must not be added to the ensemble's interval. The diagnostic that carries these quantities, including the exceedance probabilities in closed form and the surrogate comparison above, is \texttt{pakemis\_effluent\_tail\_diagnostic.csv}. It is registered in the validation catalogue with an explicit interpretation limit rather than as evidence: both of its sides read the same published range, so no external evidence enters it, and it establishes what the decomposition decomposes rather than whether the factor is right.

The estimator is Jansen pick-freeze on two independent sample matrices with \num{32768} samples per matrix, at a cost of $(k+2)N$ surrogate evaluations. The independent sampling units are dependency groups and stream-specific residuals rather than individual emission rows, so a shared parameter enters once; the residual bucket is reported as a unit of its own, so the share of variance that carries no parameter name is visible rather than appearing as the gap between the indices and one.

\textbf{Only total-effect indices are published.} First-order indices and the interaction term derived from them were withdrawn as typed absences: $\sum_i S_i \le 1$ is an exact bound, and at the published sample size the estimates summed to \num{1.325}, a third above it, with the consequence that the derived interaction was negative on \num{142} of \num{151} rows. Carrying the measurement over a sixteen-fold range of sample size shows $\sum_i S_i$ swinging from \num{0.447} to \num{1.325} to \num{1.169} while $\sum_i S_{Ti}$ holds between \num{1.021} and \num{1.030}, which locates the failure in the estimator rather than in the model. That sweep ships as \texttt{pakemis\_variance\_estimator\_sweep.csv}, so the withdrawal can be re-run against a changed sampler rather than taken on trust. What is withdrawn is an estimate and not a finding: this does not say interactions are absent, it says this estimator cannot resolve them at a sample size the build can afford. Every reported ranking reads the total-effect index, and total-effect indices include interactions and need not sum to one. The stability of $\sum_i S_{Ti}$ across the sweep is evidence about that \emph{sum} and about the first-order estimator that failed beside it; it is not evidence that the individual total-effect indices have converged, and no claim here rests on the ordering of two parameters whose indices are close. Nor does a high variance share by itself make a parameter the largest driver of the upper tail: variance importance is a statement about spread around the centre, and the tail is governed by the shape of each factor's support as well as by its variance. This analysis is conditional on the registered parameterization and is distinct from the category-level Approach~1 variance ranking.

Every external and internal check the inventory makes is registered rather than described, so a reader can ask what a given comparison can and cannot establish without reconstructing it from prose. The register holds \num{103} entries, each carrying the property tested (magnitude, timing or geography), the module, the grain, the comparator and its version, the period and boundary on each side, the inputs the two sides are known to share, an independence class, whether the entry is a validation, a corroboration, a calibration or a declared absence, the output that holds its result, and --- for all \num{103} --- an explicit interpretation limit.

The independence classes are the reason the register exists, and they partition the register exactly: of \num{103} entries, \num{11} compare against an independent compilation, \num{45} are partial, \num{20} share a large part of their inputs, \num{2} share an upstream statistic outright, \num{5} are internal identities or consistency checks, \num{12} have no applicable comparator, and \num{8} have none. A second, \emph{independent} classification records what each entry is for rather than how independent it is, and it partitions the same \num{103}: \num{70} validations, \num{15} corroborations, \num{7} calibrations and \num{11} declared absences. The two must not be added together --- a calibration also has an independence class, and a breakdown that mixed the two columns would double-count the \num{3} calibrations that sit inside the partial and large classes while omitting the entries whose class it never named. Calibrations are separated because an input fitted to a comparator cannot then be validated by it; declared absences are carried rather than omitted, so the gaps are counted alongside the checks.
Table~\ref{tab:five_classes} regroups the same \num{103} entries by what each comparison can establish, mapping the register's independence and role columns onto five classes (\texttt{campaign/v3\_readiness/validation\_catalog\_five\_classes.csv}).
%% GENERATED by scripts/make_five_classes_table.py from pakemis_validation_five_classes_summary.csv. Do not edit.
\begin{table}[t]
\caption{The validation catalogue regrouped by what each entry can establish. The count of entries is a register count, not a count of validations; the class of every entry is in \texttt{outputs/pakemis\_validation\_five\_classes.csv}.}
\label{tab:five_classes}
\footnotesize
\begin{tabular}{@{}>{\raggedright\arraybackslash}p{6.0cm}r@{}}
\tophline
Class & Entries \\
\middlehline
Partially independent benchmark & 62 \\
Independent compilation (another inventory, not an observation) & 11 \\
Calibration (fitted to the comparator; cannot validate it) & 7 \\
Shares activity or factors with the estimate & 6 \\
Internal consistency or accounting identity & 5 \\
Declared absence (no check performed) & 11 \\
Not a comparison & 1 \\
\middlehline
All catalogue entries & 103 \\
\bottomhline
\end{tabular}
\end{table}
 The class of every entry is in \texttt{outputs/pakemis\_validation\_five\_classes.csv} and the counts in \texttt{outputs/pakemis\_validation\_five\_classes\_summary.csv}; both are generated from the catalogue by one rule on its role and independence columns. The terms of every external comparison --- source, gas metric, year basis on each side, units, category boundary and independence class --- are declared row by row in \texttt{outputs/pakemis\_comparison\_crosswalk.csv}, built from the catalogue and the BTR table so it cannot drift from them.

A companion benchmark table adds, per comparison, whether the two boundaries match and what ratio the comparison is permitted to support. Together these decide what a number like a \num{0.96} national ratio means: without the boundary and shared-input fields it reads as accuracy, and with them it reads as what it is.

Structural alternatives address different interpretations of evidence, including the power coal-energy basis, sewer methane correction factor, fugitive operating practices and residential dung use. The reference configuration fixes these choices. Each alternative is evaluated through the relevant calculation and reported with its change from the reference. Alternatives are not assigned probabilities or pooled into the parameter ensemble. Mutually exclusive readings of the same quantity are identified as a family, and changes from separate alternatives are not added unless their joint calculation is explicitly defined. Unestimated categories retain their notation keys; comparator-based magnitudes and scope screenings indicate possible omissions without being inserted into the reference total.

The annual, monthly and provincial uncertainty products share realization identifiers. A 10\,000-draw subset is identical to the corresponding prefix of the national run. Monthly weights are sampled only where measured profile dispersion is available, using the documented error model and renormalizing each realization to sum to one. Elsewhere the profile remains fixed or is represented by a separate conditional alternative. Provincial shares remain fixed within a configuration; alternative allocations are evaluated on the same emission draws. Each realization satisfies
\begin{equation}
 \begin{split}
 \sum_m E_{c,m}^{(d)}&=E_{c,\mathrm{annual}}^{(d)},\\
 \sum_p E_{c,p}^{(d)}+E_{c,\mathrm{unallocated}}^{(d)}
 &=E_{c,\mathrm{annual}}^{(d)}.
 \end{split}
 \label{eq:uncertainty_allocation}
\end{equation}
Parameter intervals and temporal, spatial or completeness brackets remain separately identified. Any envelope across alternatives covers only the configurations considered and is not a probability interval or a comprehensive bound on omitted uncertainty.

\begin{table*}[t]
\caption{Estimated emissions by IPCC category and gas in FY2023--24 under the reference configuration. \textbf{Codes are IPCC 2006 throughout} (Volume~1, Chapter~8, Table~8.2); a crosswalk to the UNFCCC common reporting tables, including the categories that have no counterpart in either direction, is published as \texttt{outputs/pakemis\_category\_crosswalk.csv}, which carries beside each rule the components that resolved through it and what they sum to. An earlier edition of this table mixed the two schemes, reporting enteric fermentation at the CRT's 3.A beside solid-waste disposal at IPCC 2006's 4.A. All gas columns and totals are Mt~CO$_2$e on the AR5 basis. A dash denotes no contribution to this reported row, not a national completeness notation. Rounded cells can differ slightly from rounded row or column totals.}
\label{tab:categories}
\small
\begin{tabular}{lp{0.34\textwidth}rrrr}
\tophline
IPCC code & Category & CO$_2$ & CH$_4$ & N$_2$O & Total \\
\middlehline
3.A.1 & Enteric fermentation & -- & 151.658 & -- & 151.658 \\
1.A.3.b & Road transportation & 39.285 & 0.398 & 0.511 & 40.194 \\
1.A.1.a & Public electricity and heat & 39.907 & 0.027 & 0.106 & 40.040 \\
1.A.4.b & Residential & 18.388 & 9.835 & 1.244 & 29.467 \\
1.A.2 & Manufacturing and construction & 28.009 & 0.053 & 0.069 & 28.131 \\
4.D.1 & Domestic wastewater & -- & 21.886 & 3.034 & 24.919 \\
1.B.2 & Oil and gas fugitives & 0.916 & 23.777 & 0.001 & 24.694 \\
2.A.1 & Cement production & 21.816 & -- & -- & 21.816 \\
3.A.2 & Manure management (CH$_4$ and direct N$_2$O) & -- & 13.182 & 6.148 & 19.330 \\
3.C.7 & Rice cultivation & -- & 16.548 & -- & 16.548 \\
3.C.4 & Direct N$_2$O from managed soils & -- & -- & 16.434 & 16.434 \\
3.C.5 & Indirect N$_2$O from managed soils & -- & -- & 12.091 & 12.091 \\
4.A & Solid-waste disposal & -- & 9.169 & -- & 9.169 \\
3.C.3 & Urea application & 4.824 & -- & -- & 4.824 \\
4.C.2 & Open burning of waste & 1.853 & 1.711 & 0.262 & 3.826 \\
2.B.1 & Ammonia production & 3.243 & -- & -- & 3.243 \\
3.C.6 & Indirect N$_2$O from manure management & -- & -- & 1.898 & 1.898 \\
1.A.4.c$^{a}$ & Mixed off-road machinery & 1.569 & 0.003 & 0.160 & 1.732 \\
3.C.1.b & Biomass burning in croplands & -- & 0.538 & 0.132 & 0.670 \\
2.C.1 & Iron and steel production & 0.393 & -- & -- & 0.393 \\
1.A.3.c & Railway traction & 0.353 & 0.001 & 0.036 & 0.389 \\
1.A.3.a & Domestic aviation & 0.299 & 0.000 & 0.002 & 0.301 \\
2.D.1 & Lubricant use & 0.083 & -- & -- & 0.083 \\
$^{b}$ & Railway non-traction & 0.005 & 0.000 & 0.000 & 0.005 \\
\middlehline
 & Total & 160.944 & 248.785 & 42.129 & 451.858 \\
\bottomhline
\end{tabular}
\belowtable{$^{a}$ The mixed off-road allocation includes uses not fully separated between 1.A.4.c and 1.A.2. $^{b}$ The 1.A.2/1.A.4.c split is unresolved. Public electricity excludes captive industrial generation, which contributes to 1.A.2.}
\end{table*}

Trend uncertainty is calculated by differencing corresponding realizations over a common sector scope. The reported temporal dependency setting holds emission factors, physical constants and other persistent parameters across years, while allowing year-specific activity to vary between years. Sensitivity settings increase or decrease temporal sharing. For years $a$ and $b$, each change is calculated as $\Delta T^{(d)}=T_b^{(d)}-T_a^{(d)}$ before taking percentiles. This allows common errors to cancel and avoids treating two annual level intervals as an interval for their difference.



\section{Reference-year category and gas account}
\label{sec:category_account}

Table~\ref{tab:categories} retains the full reference-year category account. The columns report each gas's CO$_2$-equivalent contribution, not native gas mass. Modules and IPCC categories differ where one category receives contributions from several modules; public generation and captive industrial generation remain separately allocated. These are estimates for the implemented scope, not a complete national inventory.



\appendix
\section{Notation and reporting conventions}
\label{app:notation}

Table~\ref{tab:notation} summarises the notation and reporting conventions used in the inventory.

\begin{table*}[t]
\caption{Conventions needed to interpret the inventory.}
\label{tab:notation}
\small
\begin{tabular}{p{0.24\textwidth}p{0.66\textwidth}}
\tophline
Term & Meaning \\
\middlehline
FY2023--24 & 1 July 2023 to 30 June 2024; analogous boundaries apply to other fiscal years. \\
Native gas & Mass of CO$_2$, CH$_4$ or N$_2$O before GWP weighting. \\
AR5 CO$_2$e & GWP$_{100}$: CO$_2$ 1; biogenic CH$_4$ 28; fossil CH$_4$ 30; N$_2$O 265. \\
AR6 CO$_2$e & GWP$_{100}$: CO$_2$ 1; biogenic CH$_4$ 27; fossil CH$_4$ 29.8; N$_2$O 273. \\
Biogenic CO$_2$ & Biomass-combustion memo item, excluded from the reported non-LULUCF total. \\
L and S & Agricultural-diversion and household-demand interpretations of residential dung, respectively. \\
NE; IE; NO; NA & Not estimated; included elsewhere; not occurring; not applicable. These are not interchangeable with zero. \\
Unallocated & Estimated mass lacking placement at the requested spatial resolution; retained in national reconciliation. \\
Conditional interval & Parameter percentiles under a stated scope, configuration and dependence model. \\
Structural alternative & A separate interpretation of evidence or system boundaries with no assigned probability. \\
\bottomhline
\end{tabular}
\end{table*}

\noappendix

\begin{thebibliography}{99}

\bibitem[Crippa et~al.(2020)]{Crippa2020}
Crippa, M., Solazzo, E., Huang, G., Guizzardi, D., Koffi, E., Muntean, M., Schieberle, C., Friedrich, R., and Janssens-Maenhout, G.: High resolution temporal profiles in the Emissions Database for Global Atmospheric Research, Scientific Data, 7, 121, \doi{10.1038/s41597-020-0462-2}, 2020.

\bibitem[FAO(2025)]{FAOSTAT2025}
FAO: FAOSTAT: emissions totals, crops and livestock, fertilisers by nutrient, and forestry production and trade, Food and Agriculture Organization of the United Nations, Rome [data sets], \url{https://www.fao.org/faostat/}, 2025 snapshot used in the source inventory; catalogue checked 30 September 2026, 2025.

\bibitem[IPCC(2006)]{IPCC2006}
IPCC: 2006 IPCC Guidelines for National Greenhouse Gas Inventories, edited by: Eggleston, H. S., Buendia, L., Miwa, K., Ngara, T., and Tanabe, K., Institute for Global Environmental Strategies, Hayama, Japan, \url{https://www.ipcc-nggip.iges.or.jp/public/2006gl/}, 2006.

\bibitem[IPCC(2013)]{IPCCAR5}
IPCC: Climate Change 2013: The Physical Science Basis, Cambridge University Press, Cambridge, \doi{10.1017/CBO9781107415324}, 2013.

\bibitem[IPCC(2019)]{IPCC2019}
IPCC: 2019 Refinement to the 2006 IPCC Guidelines for National Greenhouse Gas Inventories, edited by: Calvo Buendia, E., Tanabe, K., Kranjc, A., Baasansuren, J., Fukuda, M., Ngarize, S., Osako, A., Pyrozhenko, Y., Shermanau, P., and Federici, S., IPCC, Switzerland, \url{https://www.ipcc-nggip.iges.or.jp/public/2019rf/}, 2019.

\bibitem[IPCC(2021)]{IPCCAR6}
IPCC: Climate Change 2021: The Physical Science Basis, Cambridge University Press, Cambridge, \doi{10.1017/9781009157896}, 2021.

\bibitem[Kaza et~al.(2018)]{Kaza2018}
Kaza, S., Yao, L., Bhada-Tata, P., and Van Woerden, F.: What a Waste 2.0: A Global Snapshot of Solid Waste Management to 2050, World Bank, Washington, DC, \doi{10.1596/978-1-4648-1329-0}, 2018.

\bibitem[Ministry of Finance(2024)]{MoF2024}
Ministry of Finance: Pakistan Economic Survey 2023--24, Government of Pakistan, Islamabad, chapter 2, Agriculture, 2024.

\bibitem[NEPRA(2024)]{NEPRA2024}
NEPRA: State of Industry Report 2024, National Electric Power Regulatory Authority, Islamabad, \url{https://nepra.org.pk/publications/State%20of%20Industry%20Reports/State%20of%20Industry%20Report%202024.pdf}, 2024.

\bibitem[NFDC(2024)]{NFDC2024}
NFDC: Fertilizer Review 2023--24, Publication No. 2/2024, National Fertilizer Development Centre, Government of Pakistan, Islamabad, \url{https://nfdc.gov.pk/publications/Annual%20Fertilizer%20Review%202023-24.pdf}, 2024.

\bibitem[Pak-EPA and EPMC(1996)]{PakEPA1996}
Pak-EPA and EPMC: National Study on Privatization of Solid Waste Management in Eight Cities of Pakistan, Pakistan Environmental Protection Agency, Islamabad, 1996.

\bibitem[PBS(2019)]{HIES2019}
PBS: Household Integrated Economic Survey 2018--19, Pakistan Bureau of Statistics, Islamabad, survey report and microdata, 2019.

\bibitem[PBS(2023)]{PBSCensus2023}
PBS: Seventh Population and Housing Census 2023, Pakistan Bureau of Statistics, Islamabad, census results, 2023.

\bibitem[PBS(2026a)]{PBSAgri2026}
PBS: Agricultural Census 2024, Pakistan Bureau of Statistics, Islamabad, report collection, \url{https://www.pbs.gov.pk/agriculture-census/}, accessed 30 September 2026, 2026a.

\bibitem[PBS(2026b)]{PBSHIES2026}
PBS: Household Integrated Economic Survey (HIES) 2024--25, Pakistan Bureau of Statistics, Islamabad, \url{https://www.pbs.gov.pk/household-integrated-economic-survey-hies-2024-25/}, accessed 30 September 2026, 2026b.

\bibitem[WHO and UNICEF(2023)]{JMP2023}
WHO and UNICEF: Joint Monitoring Programme for Water Supply, Sanitation and Hygiene: Pakistan household country file, 2023 snapshot, \url{https://washdata.org/data/downloads}, 2023.

\end{thebibliography}

\end{document}
